%قبل از شروع مسئله ترجیحاً چیزی ننویسید. اما اگر خواستید بنویسید، در قالب لیستی (فقرات و شمارش) بنویسید.
%برای اضافه کردن نمره، بعدی از شروع‌مسئله، [x] را اضافه کنید و به جای x  نمره مورد نظر را قرار دهید. برای زیرمسئله‌ها هم به همین صورت است.
%\شروع‌مسئله[10]
\شروع‌مسئله [25]
% صورت اصلی مسئله را در این‌جا بنویسید.
	\متن‌لاتین{altF4}، قهرمان داستان ما، علاقۀ بسیار زیادی به بازی کردن با آرایه‌ها دارد و تخصص اصلی او
	در مرتب کردن آرایه‌هاست. به تازگی دوستان او، \متن‌لاتین{altF4} را با یک چالش جدید روبرو کرده‌اند.
	چالش به این صورت است که در ابتدا به او یک آرایه به طول $n$از اعداد صحیح داده می‌شود و وظیفه \متن‌لاتین{altF4}
	این است که این آرایه را به صورت صعودی مرتب کند. اما مشکل اصلی این است که دوستانش یک سری شرط برای 
	این کار گذاشته‌اند. از جمله این که \متن‌لاتین{altF4} در هر حرکت می‌تواند جای دو عنصر مجاور \متن‌لاتین{a[i]}	و \متن‌لاتین{a[i+1]} را با هم عوض کند.(به ازای
	$1 \leq i \leq n-1 $،
	 اولین عنصر آرایه \متن‌لاتین{a[1]} در نظر گرفته شده است)
	دوستان \متن‌لاتین{altF4} پا را فراتر گذاشته‌اند و این شرط را هم اضافه کرده‌اند که با جابه‌جا کردن 
	\متن‌لاتین{a[i]} و \متن‌لاتین{a[i+1]}، مقدار \متن‌لاتین{a[i]} یک واحد کاهش و مقدار \متن‌لاتین{a[i+1]} 
	یک واحد افزایش می‌یابد. او عملیات جا‌به‌جایی را به هر تعداد که بخواهد، می‌تواند انجام دهد. از آنجا 
	که \متن‌لاتین{altF4} به همه گفته که تخصصش، مرتب کردن آرایه‌هاست و الان از پس این کار برنمی‌آید، به
	سراغ شما آمده تا مشکل او را برطرف کنید.
	
	\متن‌سیاه{ورودی:}
	
	در خط اول ورودی اندازه آرایه ورودی یعنی $n$ می‌آید.
	
	در خط دوم ورودی $n$ عدد که اعضای آرایه هستند با یک فاصله از هم می‌آیند.
	
	\متن‌سیاه{خروجی:}
	
	در خروجی، اگر آرایه مرتب‌شده وجود نداشت، عبارت \متن‌لاتین{No} چاپ می‌شود.
	
	اگر آرایه مرتب‌شده وجود داشت، در خط اول عبارت  \متن‌لاتین{Yes} و در خط بعدی آرایه مرتب شده چاپ می‌شود.
	
	\متن‌سیاه{محدودیت‌ها:}
	
	$1 \leq N \leq 10^5$\\
	$1 \leq a[i] \leq 10^9$\\
	\متن‌سیاه{ورودی و خروجی نمونه:}
	
%	تست کیس اول
	\begin{latin}
		\begin{tabularx}{\textwidth}{|X | X|}
			\hline
			\textbf{\rl{ورودی}} &
			\textbf{\rl{خروجی}}\\
			\hline
			4&
			Yes\\
			7 3 9 10 &
			4 6 9 10\\
			\hline
		
		\end{tabularx}
	\end{latin}

%	تست کیس دوم
		\begin{latin}
		\begin{tabularx}{\textwidth}{|X | X|}
			\hline
			\textbf{\rl{ورودی}} &
			\textbf{\rl{خروجی}}\\
			\hline
			3&
			No\\
			11 1 9 &
			\\
			\hline
			
		\end{tabularx}
	\end{latin}


%اگر مسئله چند قسمت داشت به این صورت عمل کنید:
%\شروع{زیرمسائل}
%\زیرمسئله[0] % مسئلۀ ۱آ
%\زیرمسئله % مسئلۀ ۱ب
%\زیرمسئله % مسئلۀ ۱پ
%\زیرمسئله % مسئلۀ ۱ت
%\پایان{زیرمسائل}

\solution{
	پاسخ نمونه.
}




%%%%%%%%%%%%%%%%%%%%%%%%%%%%%%%%%%%%%%%%%%%%%%%%%%%%%
% برخی دستورات مفید برای تمرین‌ها. %
%%%%%%%%%%%%%%%%%%%%%%%%%%%%%%%%%%%%%%%%%%%%%%%%%%%%%

% فرمول‌های نمونه:
%$f(x) = \Theta(x)$
%$T(x, y) \leq \log(x) + 2^y + \binom{2n}{n}$
% {\tt code\_function}


% لیست موارد بدون شماره‌گذاری:
%\شروع{فقرات}
%    \فقره مورد اول 
%        \شروع{فقرات}
%            \فقره مورد اول داخلی 
%                \شروع{فقرات}
%                    \فقره اولین مورد خیلی داخلی 
%                    \فقره دومین مورد خیلی داخلی 
%                \پایان{فقرات}
%            \فقره مورد دوم داخلی 
%        \پایان{فقرات}
%    \فقره مورد دوم 
%\پایان{فقرات}

% لیست موارد شماره‌گذاری شده:
%\شروع{شمارش}
%    \فقره یک مورد 
%    \فقره یک مورد دیگر 
%    \فقره باز هم یک مورد دیگر
%\پایان{شمارش}

% برای نوشتن یک فرمول چند خطی که یک شماره بخوره. برای حالت‌های دیگر به سایت زیر مراجعه کنید:
% https://www.overleaf.com/learn/latex/Aligning_equations_with_amsmath
%\begin{align} 
%    \begin{split}
%        (x+y)^3 &= (x+y)^2(x+y) \\
%                &= (x^2+2xy+y^2)(x+y) \\
%                &= (x^3+2x^2y+xy^2) + (x^2y+2xy^2+y^3) \\
%                &= x^3+3x^2y+3xy^2+y^3
%    \end{split}                                 
%\end{align}

% نوشتن ماتریس و آرایه دوبعدی
%\begin{align}
%    A = 
%    \begin{bmatrix}
%        A_{11} & A_{21} \\
%        A_{21} & A_{22}
%    \end{bmatrix}
%\end{align}

% اضافه کردن تصویر با شرح. برای اضافه کردن تصویر بدون شرح، فقط دستور درج‌تصویر را استفاده کنید.
%\شروع{شکل}[h!]
%\تنظیم‌ازوسط
%  \درج‌تصویر[پهنا=0.5\پهنای‌متن]{img.jpg}
%\شرح{شرح تصویر}
%\برچسب{تصویر:۱}
%\پایان{شکل}
%\پایان{زیرمسائل}

%کد زیر را ببینید.
%%\lstinputlisting{source_filename.py}
%\begin{latin}
%\begin{lstlisting}
%# Import the modules
%import sys
%import random
%
%ans = True
%
%while ans:
%question = raw_input("Ask the magic 8 ball a question: (press enter to quit) ")
%
%answers = random.randint(1,8)
%
%if question == "":
%sys.exit()
%
%elif answers == 1:
%print "It is certain"
%
%elif answers == 2:
%print "Outlook good"
%
%elif answers == 3:
%print "You may rely on it"
%
%elif answers == 4:
%print "Ask again later"
%
%elif answers == 5:
%print "Concentrate and ask again"
%
%elif answers == 6:
%print "Reply hazy, try again"
%
%elif answers == 7:
%print "My reply is no"
%
%elif answers == 8:
%print "My sources say no"
%\end{lstlisting}
%\end{latin}
